\chapter{Midterm Practice}
\section{HW6}
\begin{problem}
    Let \(A\) be an open interval in \(\mathbb{R}\). Show that 
    \[
        m^*(A) = m^*(A \cap (0, \infty)) + m^*(A \setminus (0, \infty)).
    \]  
\end{problem}

\begin{proof}
    Let \(A = (\alpha , \beta)\), then 
    \[
        A \cap (0, \infty) = \begin{dcases}
            \varnothing, &\text{ if } \beta \le 0 ;\\
            (\alpha , \beta), &\text{ if } \alpha , \beta > 0 ;\\
            (0 , \beta), &\text{ if } \alpha \le 0, \beta > 0,
        \end{dcases}
    \] 
    and 
    \[
        A \setminus (0, \infty) = \begin{dcases}
            (\alpha , \beta), &\text{ if } \beta \le 0 ;\\
            \varnothing, &\text{ if } \alpha, \beta > 0 ;\\
            (\alpha , 0], &\text{ if } \alpha \le 0, \beta > 0.
        \end{dcases}
    \]
    Thus, in each case 
    \[
        m^*(A) = m^*(A \cap (0, \infty)) + m^*(A \setminus (0, \infty)).
    \]
\end{proof}

\begin{problem}
    Let \(A\) be an open box in \(\mathbb{R} ^n\), and let \(E = \left\{ (x_1, \dots , x_n) \in \mathbb{R} ^n : x_n > 0 \right\} \). Show that 
    \[
        m^*(A) = m^*(A \cap E) + m^*(A \setminus E).
    \]  
\end{problem}

\begin{proof}
    Let \(A = \prod _{i=1}^{\infty} (a_i, b_i)\), then 
    \[
        A \cap E = \left( \prod _{i=1}^{n-1} (a_i, b_i) \right) \times ((a_n \cap b_n) \cap (0, \infty)). 
    \] 
    Now since \((a_n, b_n) \cap (0, \infty)\) is an interval,  
    \[
        m^*(A \cap E) = \mathrm{vol}(A_{n-1}) \cdot m^*(A_n \cap (0, \infty)) \text{ where } A_{n-1} = \prod _{i=1}^{n-1} (a_i, b_i), \ A_n = (a_n, b_n). 
    \]
    Similarly, 
    \[
        A \setminus E = A_{n-1} \times (A_n \setminus (0, \infty)),
    \]
    so 
    \[
        m^*(A \setminus E) = \mathrm{vol}(A_{n-1}) \cdot m^*(A_n \setminus (0, \infty)).
    \]
    Hence,
    \begin{align*}
        m^*(A) &= \mathrm{vol}(A) = \mathrm{vol}(A_{n-1})(m^*(A_n))=\mathrm{vol}(A_{n-1})(m^*(A_n \cap (0, \infty)) + m^*(A_n \setminus (0, \infty))) \\
        &= \mathrm{vol}(A_{n-1}) \cdot m^*(A_n \cap (0, \infty)) + \mathrm{vol}(A_{n-1}) \cdot m^*(A_n \setminus (0, \infty)) \\
        &= m^*(A \cap (0, \infty)) + m^*(A \setminus (0, \infty)). 
    \end{align*}
\end{proof}

\section{HW1}
\begin{problem}
    Let 
    \[
        D = \left\{ (x, y) \in \mathbb{R} ^n, x \ge 0, 0 \le y \le x^3 \right\}, 
    \]
    and \(x_0 = (0, 0)\). Define \(f: D \to \mathbb{R}\) by \(f(x, y) = 2x + \sqrt{y} \). 
    \begin{itemize}
        \item [(a)] Prove that \(f\) is differentiable at \((0, 0)\) on \(D\). 
        \item [(b)] Describe the set of all derivative \(L(x, y)\) for \(f\) on \((0, 0)\) on \(D\).       
    \end{itemize}   
\end{problem}

\section{HW2}
\begin{problem}
    Define 
    \begin{align*}
        \mathrm{GL}(n, \mathbb{R}) &= \left\{ A \text{ is a } n\times n \text{ matrix and } A \text{ is invertible}    \right\} \\
        &= \left\{ A \mid \det (A) \neq 0 \right\},  
    \end{align*}
    which is open in \(\mathbb{R} ^{n + m}\). Consider the inverse map \(g: \mathrm{GL}(n, r) \to \mathrm{GL}(n, r)\), i.e. \(g(A) = A^{-1}\) . Show that \(g^{\prime} (A): \mathbb{R} ^{n \times n} \to \mathrm{R}^{n \times n} \) is \(g^{\prime} (A)(H) = A^{-1} H A^{-1}\) for every \(H \in \mathbb{R} ^{n \times n}\). 

    \begin{remark}
        We can use the fact that if \(M \in \mathbb{R} ^{n + m}\) has \(\lVert M \rVert < 1 \), then 
        \[
            (I + M)^{-1} = I - M + M^2 - M^3 + \dots 
        \]  
        converges in \(\mathbb{R} ^{n + m}\). 
    \end{remark}      
\end{problem}

\section{HW3}

\begin{problem}
    Let \(f: S \to S\) be a map between a complete metric space \(S\). Assume there is a real sequence \(\left\{ \alpha _n \right\} \) which converges to \(0\) s.t. 
    \[
        d(f^n(x), f^n(y)) \le \alpha _n d(x, y)
    \]    
    for all \(n \ge 1\) and all \(x, y \in S\). Prove that \(f\) has a unique fixed point. Hint. Apply the fixed-point theorem to \(f^m\) for a suitable \(m\).  
\end{problem}

\section{HW4}

\begin{problem}
    Let \(n\ge 2\), and let
\[
f:\mathbf{R}^n\to \mathbf{R}
\]
be a smooth function. Define
\[
S:=\{x=(x_1,\dots,x_n)\in \mathbf{R}^n : f(x)=0\}.
\]
Assume that
\[
\nabla f(x)\neq 0
\qquad\text{for every }x\in S.
\]
Then for every point \(p\in S\), there exists an index \(j\in\{1,\dots,n\}\), an open neighborhood
\(W\subset \mathbf{R}^n\) of \(p\), an open set \(U\subset \mathbf{R}^{\,n-1}\), and a differentiable
function
\[
g:U\to \mathbf{R}
\]
such that, after possibly reordering the coordinates, the set \(S\cap W\) can be written in the form
\[
S\cap W
=
\Bigl\{
(x_1,\dots,x_{j-1},\, g(x_1,\dots,x_{j-1},x_{j+1},\dots,x_n),\, x_{j+1},\dots,x_n)
:
(x_1,\dots,x_{j-1},x_{j+1},\dots,x_n)\in U
\Bigr\}.
\]
In other words, near every point of \(S\), the hypersurface \(S\) is locally the graph of a differentiable
function of \(n-1\) variables.
\end{problem}

\section{HW5}
\begin{problem}
    Recall that a function $f: [a, b] \to \mathbf{R}$ is \textit{Lipschitz} if there exists $L \geq 0$ such that $|f(x) - f(y)| \leq L|x-y|$ for all $x, y \in [a, b]$. 
Define the graph of $f$ as:
\[
\Gamma_f := \{(x, f(x)) : x \in [a, b]\} \subset \mathbf{R}^2.
\]
Prove that $m^*_2(\Gamma_f) = 0$. 
\textit{Hint: Partition $[a, b]$ into $n$ sub-intervals and cover the corresponding pieces of the graph with rectangles. Consider the limit as $n \to \infty$.}
\end{problem}

\section{Random Stuff}
\begin{theorem}
    If \(f\) is continuoius and \(I\) is an open interval, then \(f(I)\) is also an interval.   
\end{theorem}

\begin{theorem}
    Any interval in \(\mathbb{R}\) is measurable (open, closed, or half-open, half-closed) 
\end{theorem}

\begin{theorem}
    A set \(E\) is Lebesgue measurable iff \(m^*(A) = m^*(A \cap E) + m^*(A \setminus E)\) for all \(A \subseteq \mathbb{R}^n\). To show \(E\) is measurable, it suffices to show 
    \[
        m^*(A) \ge m^*(A \cap E) + m^*(A \setminus E).
    \]    
\end{theorem}

\chapter{Additional proof}
\begin{theorem} \label{thm: linf eq lsup eq f, then converge}
    If \(f = \liminf_{n \to \infty} f_n = \limsup_{n \to \infty} f_n  \), then \(f_n \to f\) pointwise.  
\end{theorem}

\begin{proof}
Let $\{f_n\}_{n=1}^\infty$ be a sequence of real-valued functions on a set $\Omega$, and suppose that for every $x \in \Omega$,
\[
f(x) = \liminf_{n\to\infty} f_n(x) = \limsup_{n\to\infty} f_n(x).
\]

We will show that $f_n(x) \to f(x)$ for every $x \in \Omega$.

Fix $x \in \Omega$. Then $\{f_n(x)\}$ is a sequence of real numbers satisfying
\[
\liminf_{n\to\infty} f_n(x) = \limsup_{n\to\infty} f_n(x) = f(x).
\]

We prove that $f_n(x) \to f(x)$.

\medskip

Let $\varepsilon > 0$.

By the definition of $\limsup$, we have
\[
\limsup_{n\to\infty} f_n(x) = \lim_{n\to\infty} \sup_{k \ge n} f_k(x) = f(x).
\]
Hence there exists $N_1 \in \mathbb{N}$ such that for all $n \ge N_1$,
\[
\sup_{k \ge n} f_k(x) < f(x) + \varepsilon.
\]
In particular, for all $k \ge n \ge N_1$,
\[
f_k(x) \le \sup_{j \ge n} f_j(x) < f(x) + \varepsilon.
\]
Thus for all $k \ge N_1$,
\[
f_k(x) < f(x) + \varepsilon.
\]

\medskip

Similarly, by the definition of $\liminf$, we have
\[
\liminf_{n\to\infty} f_n(x) = \lim_{n\to\infty} \inf_{k \ge n} f_k(x) = f(x).
\]
Hence there exists $N_2 \in \mathbb{N}$ such that for all $n \ge N_2$,
\[
\inf_{k \ge n} f_k(x) > f(x) - \varepsilon.
\]
In particular, for all $k \ge n \ge N_2$,
\[
f_k(x) \ge \inf_{j \ge n} f_j(x) > f(x) - \varepsilon.
\]
Thus for all $k \ge N_2$,
\[
f_k(x) > f(x) - \varepsilon.
\]

\medskip

Let $N = \max\{N_1, N_2\}$. Then for all $n \ge N$,
\[
f(x) - \varepsilon < f_n(x) < f(x) + \varepsilon,
\]
which implies
\[
|f_n(x) - f(x)| < \varepsilon.
\]

\medskip

Since $\varepsilon > 0$ is arbitrary, we conclude that
\[
f_n(x) \to f(x).
\]

Since $x \in \Omega$ was arbitrary, the convergence is pointwise on $\Omega$.
\end{proof}

\begin{theorem} \label{thm: almost equal to measurable func then measurable}
Let $f:\mathbb{R}^n \to \mathbb{R}$ be a Lebesgue measurable function, and let $g:\mathbb{R}^n \to \mathbb{R}$ be a function such that $f=g$ almost everywhere. Then $g$ is also Lebesgue measurable.
\end{theorem}

\begin{proof}
Since $f=g$ almost everywhere, there exists a measurable set $N \subset \mathbb{R}^n$ such that $m(N)=0$ and
\[
f(x)=g(x) \quad \text{for all } x \in \mathbb{R}^n \setminus N.
\]

To prove that $g$ is measurable, it suffices to show that for every $a \in \mathbb{R}$, the set
\[
E := \{x \in \mathbb{R}^n : g(x) > a\}
\]
is measurable.

Fix $a \in \mathbb{R}$ and define
\[
F := \{x \in \mathbb{R}^n : f(x) > a\}.
\]
Since $f$ is measurable, the set $F$ is measurable.

We claim that
\[
E \triangle F \subseteq N,
\]
where $E \triangle F = (E \setminus F) \cup (F \setminus E)$ denotes the symmetric difference.

Indeed, let $x \notin N$. Then $f(x)=g(x)$, and hence
\[
g(x) > a \iff f(x) > a.
\]
Thus $x \in E$ if and only if $x \in F$, which implies $x \notin E \triangle F$. Therefore $E \triangle F \subseteq N$.

Since $N$ has measure zero, it is measurable. It follows that $E \triangle F$ is measurable. As $F$ is measurable and the collection of measurable sets is closed under symmetric difference, we conclude that
\[
E = F \triangle (E \triangle F)
\]
is measurable.

Since $a \in \mathbb{R}$ was arbitrary, this shows that $\{x : g(x) > a\}$ is measurable for all $a$, and hence $g$ is measurable.
\end{proof}